\documentclass{article}
\usepackage{booktabs}   
\usepackage{geometry}
\usepackage{graphicx}
\usepackage{placeins}
\geometry{margin=1in}

\title{Coding Logbook to Macroeconomics 3 Project}
\author{Philipp Cremer, Jannis Stagge, Felix Westerhoff}

\begin{document}
\maketitle

Step 1 - Replicate EHB

Replicate Table 1 and 2 of the paper using EWN and Stock Market Capitalisation Data from WDI
Results are closely aligned with values reported in tables 1 and 2

Discovered first data problems with WDI coverage after 2013
Investigated potential other sources without success

Discussed first Euro Area Extensions

Step 2 - Replicate Risk Sharing Regression

Data Process:
first try - OECD Data non PPP 
We tried to replicate the Consumption, GNI and GDP data using today's OECD National Accounts and exchange rates. 
We realized that \textcite{main} only weakly documented how they approached this topic, forcing us to make various 
individual choices on how to adjust the data. After a handful of rather unsuccessful tries, we decided to use OECD chained linked
volumes 2020 rebased since they are already PPP adjusted by the OECD.
Further, we tried to run the main regressions using PPP adjusted data for Consumption, GNI and GDP from the world bank. 
After a deeper analysis of this data, we decided against this data source, as GNI data for the 1990th is very limited. 
second - OECD chained linked volumes 2020 rebased
third - WDI data PPP

all yield negative risk sharing trends

fourth - OECD data non PPP, PPP adjustment with base year ???
fifth - old 2005 OECD data scraped from PDF with Claude Code

both give increasing trends with smaller slopes than original paper
smoothing\
We used the old 2005 OECD data scraped from the PDF to replicate the main regressions and confirm that our code is 
working, concluding that the problems before were related to data problems and not to a problem in our code. 


Step 3 - Refine EHB data

more detailed analysis of EHB data
exclude some irregular values (negative EHB or above 1 EHB) and investigate potential causes (WDI data mismatch, restructuring of Stock Exchanges, Ireland as financial hub)

Step 4 - Replication Panel Regressions

quite far off with initial data versions also yielding opposite signs and different significance patterns
We disovered the two stage approach (feasible GLS) which is only mentioned in a footnote. 
Once we managed to rebuild \textcite{main} data with OECD data, we replicated their regressions and recieved similar direction of effects, 
similar magnitudes and paper consistent significance patterns with data version 4 and 5. 


Step 5 - Extension of Home Bias data

investiage coverage of EWN, WDI data and subsequently EHB measure for longer time horizon and broader country sample
restrict to top 50 economies (US$ GDP) and EU
investigate irregularities in the data similar process as with original replication sample

Step 6 - Euro Area Home Bias

investigate Euro Area Home Bias as cross country averages and EA aggregate
comparison to US, and non EA sample average for different samples
investigation of negative EHB after 2014 (most likely due to countries dropping out of WDI but not out of EWN)

Step 7 - Risk sharing sample Extensions

extend for time
extend for sample
look at within Eurozone 
look at non-Eurozone

Step 8 - Panel Regression Extension

extend for time
extend for sample
look at Eurozone as an aggregate observation
investigate structural breaks

Step 9 - EHB Eurozone Extension

literature research for explanations on EHB patterns in EA
ECB working paper replication data download
investigation of paper and data mainly graphically

Step 10 - Risk Sharing Eurozone extension

seminal paper on risk sharing channels methodology
replicate main results of risk sharing channels for EU from working paper
apply methodology to Euro Zone
